\documentclass[../../../include/open-logic-section]{subfiles}

\begin{document}

\olfileid{his}{set}{infinitesimals}
\olsection{Infinitesimal lan Diferensiasi}

Newton lan Leibniz nemokaké kalkulus kanthi mandiri
ing pungkasan abad kaping pitulas. Salah siji
panganggoné kalkulus kang wigati banget yaiku
\emph{diferensiasi}. Kanthi kasar, diferensiasi
ngupaya menehi teges kanggo ``laju owah-owahan'',
utawa gradien, sawijining fungsi ing sawijining titik.

Iki cara kang cetha kanggo nggambaraké gagasan iku.
Delengen fungsi $f(x) = \nicefrac{x^2}{4} + \nicefrac{1}{2}$,
kang digambar ireng ing ngisor iki:
\begin{center}
	\begin{tikzpicture}[scale=1]
	\draw[->, gray] (-1,0) -- (4.25,0) node[right] {$x$};
	\draw[->, gray] (0,-0.25) -- (0,5.25) node[above] {$f(x)$};
	
	\foreach \x/\xtext in {1/1, 2/2, 3/3, 4/4}
	\draw[shift={(\x,0)}, gray] (0pt,2pt) -- (0pt,-2pt) node[below] {$\xtext$};
	
	\foreach \y/\ytext in {1/1, 2/2, 3/3, 4/4, 5/5}
	\draw[shift={(0,\y)}, gray] (2pt,0pt) -- (-2pt,0pt) node[left] {$\ytext$};
	
	\draw[oldiagcolorC] (0.5, 9/16) -- (3.5, 57/16) -- (3.5, 9/16)--cycle;
	\draw[oldiagcolorD] (.5, 9/16) -- (2.5, 33/16) -- (2.5, 9/16) -- cycle;
	\draw[oldiagcolorE] (.5, 9/16) -- (1.5, 17/16) -- (1.5, 9/16) -- cycle;
	\draw[thick] (-1,.75) parabola bend (0,0.5) (4,4.5);
	\end{tikzpicture}
\end{center}
Anggepen kita arep nemtokaké gradien fungsi ing $c =
\nicefrac{1}{2}$. Wiwitané, gambaren segitiga kang
sisi miringé menehi taksiran gradien ing titik iku,
upamané segitiga abang ing ndhuwur. Yèn $\beta$
iku dawa dhasar segitiga, dhuwuré yaiku
$f(\nicefrac{1}{2}+\beta) - f(\nicefrac{1}{2})$,
dadi gradien sisi miringé yaiku:
\[
\frac{f(\nicefrac{1}{2}+\beta) - f(\nicefrac{1}{2})}{\beta}.
\]
Dadi gradien segitiga !!{colorC} kang dhasaré
dawané~$3$ iku persis~$1$. Sisi miring segitiga
kang luwih cilik, yaiku segitiga !!{colorD}
kanthi dhasar dawa~$2$, menehi taksiran
kang luwih becik; gradiené $\nicefrac{3}{4}$.
Segitiga kang luwih cilik manèh, yaiku
segitiga !!{colorE} kanthi dhasar dawa~$1$,
menehi taksiran kang luwih becik manèh;
gradiené $\nicefrac{1}{2}$.

Segitiga kang saya cilik menehi taksiran
kang saya becik. Mula kita bisa kandha:
sisi miring segitiga kang dawa dhasaré
\emph{infinitesimal} menehi gradien ing $c =
\nicefrac{1}{2}$ iku dhéwé. Kanthi cara iki
kita olèh rumus turunan (kapisan) fungsi $f$
ing titik $c$:
\[
{f'}(c) = \frac{f(c+\beta) - f(c)}{\beta} \text{ kanthi $\beta$ infinitesimal.}
\]
Lan, kanthi kasar, iki kang diandharaké
Newton lan Leibniz.

Nanging, amarga padha ngandharaké mangkono,
kita kudu takon: apa kang diarani
\emph{infinitesimal}? Ana dilema kang abot.
Yèn $\beta = 0$, mula $f'$ ora
didhéfinisikaké kanthi becik, amarga
nglibataké pambagian nganggo $0$.
Nanging yèn $\beta > 0$, kita mung
olèh \emph{taksiran} gradien, dudu
gradiené iku dhéwé.

Iki dudu keprihatinan kang mung muncul
ing jaman sabanjuré. Mangkéné kritik
Berkeley marang para pandhèrèké Newton:
\begin{quote}
Aku ngakoni yèn pralambang bisa dienggo kanggo nandhani
samubarang utawa ora nandhani apa-apa: mula ing notasi
wiwitan $c + \beta$, $\beta$ bisa ateges tambahan
utawa ora ateges apa-apa. Nanging, teges endi waé
kang panjenengan pilih, panalaran panjenengan kudu
ajeg manut teges iku, lan aja nerusaké kanthi
nganggo rong teges sekaligus: tumindak mangkono
iku kesesatan nalar kang cetha.
(\citealt[\S{}XIII]{Berkeley1734},
peubah diganti supaya cocog karo tèks sadurungé)
\end{quote}
Kanggo mbéla kalkulus infinitesimal saka kritik
Berkeley, ana kang bisa mangsuli yèn omongan
ngenani ``infinitesimal'' mung kiasan. Wong bisa
kandha yèn angger segitigané \emph{cilik banget},
kita bakal olèh taksiran garis singgung kang
\emph{cukup becik}. Berkeley uga nduwèni
wangsulan: sanajan iku cukup kanggo rékayasa,
bab iki nglemahaké \emph{kalungguhan}
matématika, amarga
\begin{quote}
kita dikandhani yèn \emph{in rebus mathematicis errores qu\`{a}m minimi
non sunt contemnendi}. [Ing matématika,
kaluputan kang paling cilik waé ora kena
dianggep entheng.] \citep[\S{}IX]{Berkeley1734}
\end{quote}
Pethikan kang dicithak miring iku meh persis
kutipan saka karya Newton dhéwé,
\emph{Quadrature of Curves} (1704).

Bantahan filosofis Berkeley iku landhep banget.
Nanging kalkulus kabukti kasil banget, lan
para matématikawan terus migunakaké
tanpa tiba ing kaluputan.

\end{document}
